% Part B (AI for Authors): after writing. 11 live frames, plus 3 optional frames commented out at the end.
% AI reviews and reports; it does not rewrite the author's text. Examples are neutral and invented.

%%%%%%%%%%%%%%%%%%%%%%%%%%%%%%%%%%%%%%%%%%%%%%%%%%%%%%%%%%%
\begin{frame}[fragile]\frametitle{After Writing: AI as Reviewer, Not Writer}
\begin{itemize}
\item Once the draft is yours and finished, AI is useful as a tireless first reader.
\item It reports problems. You decide which are real, and you fix them in your own words.
\item Three jobs: editorial review, layout review, and the emails that follow.
\item Do not paste text you are not allowed to share. Follow the confidentiality rules in AI Basics and Responsible Use.
\end{itemize}
\end{frame}

%%%%%%%%%%%%%%%%%%%%%%%%%%%%%%%%%%%%%%%%%%%%%%%%%%%%%%%%%%%
\begin{frame}[fragile]\frametitle{CRAFT Example: Thorough Editorial Review}
\begin{itemize}
\item Review one chapter at a time. Ask for a list of issues, not a rewrite.
\end{itemize}
\begin{lstlisting}[language=HTML]
You are a strict copy editor (R). This is Chapter 4
of a book for first-time managers: [CHAPTER TEXT] (C).

List every inconsistency in terms, numbers and
spelling; every unclear sentence; and every claim
that needs a source (A). Do not rewrite anything.
Table: Page or paragraph | Issue | Why (T).
\end{lstlisting}
\end{frame}

%%%%%%%%%%%%%%%%%%%%%%%%%%%%%%%%%%%%%%%%%%%%%%%%%%%%%%%%%%%
% [ai-for-x: added 2026-11-08 | from the author-AI transcripts | continuity check]
\begin{frame}[fragile]\frametitle{CRAFT Example: Continuity Check}
\begin{itemize}
\item Ask for each conflict and where it occurs, not a rewrite.
\item You fix each one yourself, then run the check again.
\end{itemize}
\begin{lstlisting}[language=HTML]
You are a continuity editor (R). This is my full
chapter: [CHAPTER TEXT] (C).

List every conflict in the timeline, in names, ages
and numbers, and in what each person knows at that
point (A). Do not rewrite anything.
Table: Place | Conflict | Earlier mention (T).
\end{lstlisting}
\end{frame}

%%%%%%%%%%%%%%%%%%%%%%%%%%%%%%%%%%%%%%%%%%%%%%%%%%%%%%%%%%%
% [ai-for-x: added 2026-11-08 | factual checks of the author's own claims]
\begin{frame}[fragile]\frametitle{CRAFT Example: Fact-Check Your Claims}
\begin{itemize}
\item AI lists the claims and says how to check them. You open the sources and confirm.
\item Keep no number, date or quote that you have not verified yourself.
\end{itemize}
\begin{lstlisting}[language=HTML]
You are a fact checker (R). This is my chapter:
[CHAPTER TEXT] (C).

List every claim that states a number, a date, a
name or a quote (A). For each, say what kind of
source would confirm it. Do not correct me (T).
Table: Claim | Source type to check (T).
\end{lstlisting}
\end{frame}

%%%%%%%%%%%%%%%%%%%%%%%%%%%%%%%%%%%%%%%%%%%%%%%%%%%%%%%%%%%
% [ai-for-x: added 2026-11-08 | from the author-AI transcripts | patterns, and AI miscounts]
\begin{frame}[fragile]\frametitle{Read the Patterns, Verify the Numbers}
\begin{itemize}
\item AI is good at patterns: repeated words, long stretches without a break, how often the viewpoint changes.
\item It is weak at counting. Ask for the pattern, then count with your word processor or a script.
\item A pattern is a question to ask yourself, not a verdict.
\end{itemize}
\begin{lstlisting}[language=HTML]
You are an analyst of writing patterns (R): [CHAPTER TEXT] (C).
Which words repeat most, which paragraphs are the
longest, and where does the viewpoint change (A)?
List only; do not rewrite (T).
\end{lstlisting}
\end{frame}

%%%%%%%%%%%%%%%%%%%%%%%%%%%%%%%%%%%%%%%%%%%%%%%%%%%%%%%%%%%
\begin{frame}[fragile]\frametitle{CRAFT Example: Read It Back as a Reader}
\begin{itemize}
\item If the summary does not match what you meant, the passage is unclear.
\end{itemize}
\begin{lstlisting}[language=HTML]
You are a busy first-time manager reading this
chapter on a train (R): [CHAPTER TEXT] (C).

Summarise the chapter's main point in three lines,
then list the two places where you got confused
or stopped reading (A). Plain language (T).
\end{lstlisting}
\end{frame}

%%%%%%%%%%%%%%%%%%%%%%%%%%%%%%%%%%%%%%%%%%%%%%%%%%%%%%%%%%%
% [ai-for-x: added 2026-11-08 | from the author-AI transcripts | AI flatters; ask for the weak points]
\begin{frame}[fragile]\frametitle{Beware Flattery}
\begin{itemize}
\item AI tends to agree with you and to praise your work. Praise from a tool is not a review.
\item Ask for the weak points first, and forbid praise in the prompt.
\item A human editor sees what a tool does not. Use AI to prepare for one, not to replace one.
\end{itemize}
\begin{lstlisting}[language=HTML]
You are a tough but fair editor (R): [CHAPTER TEXT] (C).
List the five weakest points, most serious first.
No praise and no summary (A). One sentence each,
with the paragraph where it occurs (T).
\end{lstlisting}
\end{frame}

%%%%%%%%%%%%%%%%%%%%%%%%%%%%%%%%%%%%%%%%%%%%%%%%%%%%%%%%%%%
\begin{frame}[fragile]\frametitle{Layout Review with LaTeX}
\begin{itemize}
\item If you write in LaTeX, compile with your publisher's or your own template and look at the pages.
\item Check: overfull boxes, figures off the page, widows and orphans, broken cross-references (shown as question marks in the PDF), the table of contents.
\item AI helps read the compiler log: paste the warnings, ask what each one means and what to change.
\item You make the change and compile again. A compile never lies; an AI explanation might.
\end{itemize}
\begin{lstlisting}[language=HTML]
You are a LaTeX expert (R). Explain each warning in
this log in one line and say which matter for a
printed book (A): [LOG LINES] (C).
\end{lstlisting}
\end{frame}

%%%%%%%%%%%%%%%%%%%%%%%%%%%%%%%%%%%%%%%%%%%%%%%%%%%%%%%%%%%
\begin{frame}[fragile]\frametitle{CRAFT Example: The Proposal Email}
\begin{itemize}
\item Emails are communication, not authorship. You supply the facts; AI drafts the polite shell.
\item Read every line before you send. Check the publisher's own submission guidelines first.
\end{itemize}
\begin{lstlisting}[language=HTML]
You are an experienced author writing to an
acquisitions editor (R). My book: a practical guide
for first-time managers, about 250 pages, complete
draft ready (C).

Write a 150-word proposal email (A) that states
the audience, the one-line pitch and what is
attached. Polite, no flattery, no exclamation
marks (T).
\end{lstlisting}
\end{frame}

%%%%%%%%%%%%%%%%%%%%%%%%%%%%%%%%%%%%%%%%%%%%%%%%%%%%%%%%%%%
\begin{frame}[fragile]\frametitle{CRAFT Example: Follow-Up and Launch Post}
\begin{itemize}
\item Same pattern for a follow-up after two weeks, a reply to reviewer comments, or a launch post.
\end{itemize}
\begin{lstlisting}[language=HTML]
You are a courteous author (R). I sent a proposal
two weeks ago and have had no reply (C).

Write a 70-word follow-up email (A). Tone: warm,
brief, no pressure; end with one clear question (T).
\end{lstlisting}
\end{frame}

%%%%%%%%%%%%%%%%%%%%%%%%%%%%%%%%%%%%%%%%%%%%%%%%%%%%%%%%%%%
% [ai-for-x: added 2026-11-08 | from the author-AI transcripts | marketing copy is outside the manuscript]
\begin{frame}[fragile]\frametitle{CRAFT Example: Description, Hooks and Subject Lines}
\begin{itemize}
\item Marketing copy sits outside your manuscript. You supply the facts, AI offers options, you choose.
\item End the prompt with ``ask me questions first''. The answers sharpen the copy.
\end{itemize}
\begin{lstlisting}[language=HTML]
You are a bookshop copywriter (R). My book: a guide
for first-time managers on running meetings; the
reader's problem is meetings that go nowhere (C).

Write three 60-word descriptions and five one-line
hooks (A). Plain tone, no hype, no claim I cannot
support (T). Ask me any questions first.
\end{lstlisting}
\end{frame}

%%%%%%%%%%%%%%%%%%%%%%%%%%%%%%%%%%%%%%%%%%%%%%%%%%%%%%%%%%%
% [ai-for-x: optional 2026-11-08 | from the author-AI transcripts | commented out by default; remove the % to use]
% \begin{frame}[fragile]\frametitle{Translation and Audio: Gray Areas}
% \begin{itemize}
% \item Translating or narrating your finished book is not the same as writing it. Decide, and disclose.
% \item A common middle path: machine translation first, then a native speaker edits every page.
% \item Voice and translation tools have their own terms of use. Read them before you upload your book.
% \item Budget matters to many authors. Make your own choice, and do not judge the choices of others.
% \end{itemize}
% \end{frame}
%
% %%%%%%%%%%%%%%%%%%%%%%%%%%%%%%%%%%%%%%%%%%%%%%%%%%%%%%%%%%%
% [ai-for-x: optional 2026-11-08 | from the author-AI transcripts | commented out by default; the press-release point is unverified]
% \begin{frame}[fragile]\frametitle{Being Found When Readers Ask AI}
% \begin{itemize}
% \item Readers now ask AI tools to suggest books. A tool can only suggest what it can find and describe.
% \item Say clearly, on your website and retail page, who the book is for and what problem it solves.
% \item Use the words a reader would type, and keep the wording the same everywhere.
% \item Some commentators suggest a press release for each new book. The effect is unproven.
% \end{itemize}
% \end{frame}
%
% %%%%%%%%%%%%%%%%%%%%%%%%%%%%%%%%%%%%%%%%%%%%%%%%%%%%%%%%%%%
% [ai-for-x: optional 2026-11-08 | from the author-AI transcripts | commented out by default]
% \begin{frame}[fragile]\frametitle{Admin and Numbers}
% \begin{itemize}
% \item AI can read a sales report or a business plan and ask you questions. It has no feelings about the numbers.
% \item It can miscount. Check the totals yourself, and never decide on one reading.
% \end{itemize}
% \begin{lstlisting}[language=HTML]
% You are a business adviser for self-published authors (R).
% Here is my sales report for 12 months: [REPORT] (C).
% Which titles and months stand out (A)? Ask me
% questions before you advise (T).
% \end{lstlisting}
% \end{frame}
